\documentclass[11pt,twocolumn]{article}

\usepackage[a4paper,margin=0.75in,columnsep=0.3in]{geometry}
\usepackage[T1]{fontenc}
\usepackage{amsmath,amssymb,amsthm,mathtools}
\usepackage{bm}
\usepackage{graphicx}
\usepackage{booktabs}
\usepackage{array}
\usepackage{enumitem}
\usepackage{float}
\usepackage[numbers,sort&compress]{natbib}
\usepackage{hyperref}
\hypersetup{colorlinks=true,linkcolor=blue,citecolor=blue,urlcolor=blue}

\newtheorem{theorem}{Theorem}
\newtheorem{definition}{Definition}
\newtheorem{postulate}{Postulate}
\newtheorem{corollary}{Corollary}

\title{Unified Dynamics: Emergent Geometry, Gauge Structure, and Matter from a Tetrahedral Transport--Closure Substrate}
\author{%
Benjamin Walker Mayes\thanks{*Contact author: ben.w.mayes@gmail.com}\\[4pt]
{\small Unaffiliated, Cedar Park, Texas 78613, USA}\\[8pt]
Arthur Eugene Mayes\\
{\small Loyola University School of Law, USA}\\[4pt]
Dustin Huff\\
{\small Unaffiliated, USA}\\[4pt]
Brayden Ross Sanders\\
{\small Unaffiliated, USA}\\[4pt]
Jeffrey A.\ Dewey\\
{\small Unaffiliated, USA}\\[4pt]
Brandon Bevin\\
{\small Georgia Southern University, USA}\\[4pt]
A.\ Christopher Beckingham\\
{\small Queen's University, Canada}\\[4pt]
Jonathan Engle\\
{\small Unaffiliated, USA}%
}
\date{}

\begin{document}
\maketitle

\begin{abstract}
Building on earlier formulations of the packet-substrate program, we consolidate and extend Unified Dynamics (UD) into a framework in which spacetime geometry, gauge interactions, and matter arise from admissible transport on a tetrahedral packet substrate.
The theory is built on a single iterated function system (IFS) with contraction ratio $r=4^{-1/3}$ acting on the complete graph $K_4$.
All structural constants are derived from $K_4$ combinatorics: the overlap deficit $\Delta=7$, the state-space dimension $N=192$, the fine-structure coupling $\alpha_0=7/960$, the neutral seed threshold $S_0=43/7$, and the heavy-lift threshold $H=101/6$.
The gauge group $\mathrm{SU}(3)\times\mathrm{SU}(2)\times\mathrm{U}(1)$ emerges from $S_4$ representation theory with $\sin^2\theta_W=3/11$ at unification.
Electroweak symmetry breaking is driven by the $\mathbb{Z}_2$ ferromagnetic vacuum.
The Yang--Mills mass gap $\Delta_{\mathrm{mass}}\geq 3/\pi>0$ is proved.
A nonlinear fermion equation and quantum entanglement are derived from the substrate.
The Born rule emerges from carrier-viability selection.
The continuum Einstein--Hilbert action is the unique Banach fixed point of the IFS coarse-graining.
Twelve zero-parameter predictions are presented (Table~\ref{tab:predictions}), including $n_s=185/192$ and $m_\tau/m_\mu=101/6$.
The theory has one input scale ($\nu_0\approx 83$~keV) and zero free parameters.
\end{abstract}

%═══════════════════════════════════════════════════════════
\section{Introduction}
\label{sec:intro}
%═══════════════════════════════════════════════════════════

A persistent challenge in theoretical physics is the simultaneous derivation of quantum matter, gauge interactions, and gravity from a single ontology.
General relativity~\cite{einstein1915} treats geometry as dynamical but classical; quantum field theory assumes a fixed background arena.
Unified Dynamics (UD) proposes a different starting point: physical reality is a discrete packet-routing substrate whose admissible transport, closure, and selective retention generate geometry, gauge structure, and matter as derived consequences.

The defect-based substrate program was introduced in Ref.~\cite{mayes2026dd}, where localized routing defects were identified as matter excitations and internal routing symmetries as gauge sectors.
The gravitational branch was isolated in Ref.~\cite{mayes2026grav}, which established a four-channel accommodation architecture and demonstrated that Einstein dynamics emerges as the relaxed coarse isotropic limit.
The present paper extends this program by deriving all structural constants from the combinatorics of $K_4$, proving a Yang--Mills mass gap, closing the electroweak symmetry-breaking mechanism, deriving a nonlinear fermion equation, and establishing the quantum measurement chain without a measurement axiom.

The framework is organized around a single primitive: the tetrahedron $K_4$, the unique self-dual minimal 3D closure.
An iterated function system (IFS)~\cite{hutchinson1981,falconer2003} with four contractions of ratio $r=4^{-1/3}$ generates a fractal lattice whose combinatorial constants---overlap deficit $\Delta=7$, state-space dimension $N=192$---determine all physical constants without fitting.

The present paper consolidates the derivation chain from $K_4$ through gauge structure, symmetry breaking, matter, entanglement, and cosmology, with explicit equations and zero-parameter predictions.
We separate native UD derivations from continuum bridges following the translation discipline: derive natively, translate carefully, speculate visibly.

%═══════════════════════════════════════════════════════════
\section{Foundational Axioms}
\label{sec:axioms}
%═══════════════════════════════════════════════════════════

\begin{postulate}[Packet substrate]
Reality is modeled as a discrete network of packet units on the tetrahedral IFS lattice.
Each site $i$ carries state variables
\begin{equation}
X_i = (\rho_i, \theta_i, \chi_i, \sigma_i, b_i),
\label{eq:packet}
\end{equation}
where $\rho_i$ is burden density, $\theta_i\in[0,2\pi)$ is transport phase, $\chi_i\in\{0,1\}$ is closure state, $\sigma_i\in\{+1,-1\}$ is framing, and $b_i\in[0,1]$ is boundary integrity.
\end{postulate}

\begin{postulate}[Transport]
Neighboring packets are connected by transport operators $U_{ij}$ whose composition around a closed loop $C$ defines the holonomy
\begin{equation}
U_C = \prod_{(ij)\in C} U_{ij}.
\label{eq:holonomy}
\end{equation}
Nontrivial holonomy ($U_C\neq I$) signals curvature~\cite{wilson1974}.
\end{postulate}

\begin{postulate}[Admissibility]
Not every transport configuration persists.
The physical state space is
$\mathcal{T}_{\mathrm{phys}}=\{\Psi\in\mathcal{T}\mid\mathcal{A}_{\mathrm{closure}}(\Psi)=0\}$.
\end{postulate}

Six axioms further constrain the grammar:
\textbf{(A1)~Seed}---every persistent family begins from a minimal closure state;
\textbf{(A2)~Threshold}---new regimes require threshold operators costing at least~$S_0$;
\textbf{(A3)~Lift}---higher structures are generated from seeds by iterated transport;
\textbf{(A4)~Orientation}---the $\mathbb{Z}_2$ framing $\sigma_i$ determines statistics;
\textbf{(A5)~Exchange}---transport modes between closures define interactions;
\textbf{(A6)~Dressing}---observables are bare spines plus regime-conditioned corrections.

%═══════════════════════════════════════════════════════════
\section{Tetrahedral Combinatorics and IFS}
\label{sec:combinatorics}
%═══════════════════════════════════════════════════════════

The tetrahedron $K_4$ has $V=4$ vertices, $E=6$ edges, $F=4$ faces, Euler characteristic $\chi=2$, and symmetry group $S_4$ ($|S_4|=24$).

The IFS consists of four contraction maps $\{f_k\}_{k=1}^4$, each centered at a vertex of $K_4$ with ratio $r=4^{-1/3}$.
At level $n$ the lattice has $4^n$ copies with edge length $a_n=a_0 r^n$, $a_0=\sqrt{3/4}$.
The unique attractor~\cite{hutchinson1981} is the Sierpi\'nski tetrahedron.

The state-space dimension, overlap deficit, master scalar, and coherence retention are
\begin{align}
N &= |S_4|\times 2^3 = 192, \label{eq:N}\\
\Delta &= \tbinom{V}{2}+1 = 7, \label{eq:Delta}\\
x &= \Delta/N = 7/192, \label{eq:x}\\
\eta_R &= 1-x = 185/192. \label{eq:eta}
\end{align}

%═══════════════════════════════════════════════════════════
\section{Packet Action}
\label{sec:action}
%═══════════════════════════════════════════════════════════

The packet action on the IFS lattice is
\begin{multline}
S_{\mathrm{pkt}} = \sum_C \kappa_C\,\mathrm{Tr}(I{-}U_C) + \sum_i \lambda_B(\rho_i{-}\rho_0)^2 \\
+ \sum_{\langle ij\rangle}\eta_O(\theta_i{-}\theta_j)^2,
\label{eq:Spkt}
\end{multline}
where $\kappa_C=3/(4\pi)$, $\lambda_B=\kappa_C S_0$, $\eta_O=1/2$.

%═══════════════════════════════════════════════════════════
\section{Transport Seed and Obstruction}
\label{sec:transport}
%═══════════════════════════════════════════════════════════

The transport seed factorizes as
\begin{equation}
T_0 = \tfrac{1}{2}\cdot\tfrac{7}{16}\cdot\tfrac{1}{6} = \frac{7}{192},
\label{eq:seed}
\end{equation}
with gates $\eta_O^0=1/2$ (overlap), $\eta_B^0=7/16$ (branch viability), $\eta_C^0=1/6$ (coherence band).
Local transport deforms smoothly:
\begin{equation}
T(\tau) = \tfrac{7}{192}\,e^{-(\sigma_O+\sigma_B+\sigma_C)}.
\label{eq:deformed}
\end{equation}
The path action selects the least-strain coherent path:
\begin{equation}
S(\gamma) = \int_\gamma \!\Big[\log\!\tfrac{192}{7} + \sigma_O{+}\sigma_B{+}\sigma_C\Big]ds.
\label{eq:pathaction}
\end{equation}

%═══════════════════════════════════════════════════════════
\section{Derived Structural Constants}
\label{sec:constants}
%═══════════════════════════════════════════════════════════

\begin{theorem}[Neutral seed threshold]
\label{thm:S0}
$S_0 = (E\Delta+1)/\Delta = 43/7$.
\end{theorem}
\textit{Proof sketch.}
A closed loop covers all $E=6$ edges (cost~$E$); holonomy closure over $\Delta=7$ deficit modes adds the Euler term $+1$; normalizing gives $(42+1)/7$.

\begin{theorem}[Heavy-lift threshold]
\label{thm:H}
$H = 2\cdot 2^{\deg(v)}+1-\eta_C^0 = 101/6$.
\end{theorem}
\textit{Proof sketch.}
Local state space $2^3=8$; particle--antiparticle doubling gives $16$; orientation selection $+1$; coherence rebate $-1/6$.

\begin{theorem}[Interface dimension]
\label{thm:alpha}
$\alpha_0 = \Delta/(N\cdot 5)=7/960$, where $5=\dim(\mathrm{Sym}^2_0(\mathbb{R}^3))$.
\end{theorem}
\textit{Proof sketch.}
Equal coherence across all traceless symmetric deformation modes ($3\cdot 4/2-1=5$) divides the transport cost: $\alpha_0=x/5$.

The curvature coupling from discrete Gauss--Bonnet~\cite{hirani2003}:
\begin{equation}
\kappa_C = \frac{3}{4\pi}.
\label{eq:kappaC}
\end{equation}
All constants appear in Table~\ref{tab:constants}.

\begin{table}[t]
\centering
\caption{Derived constants from $K_4$.
One input $\nu_0\approx 83$~keV; zero fitted.}
\label{tab:constants}
\small
\begin{tabular}{@{}l@{\hskip 6pt}l@{\hskip 6pt}l@{\hskip 6pt}l@{}}
\toprule
Const.\ & Value & Source & Acc.\ \\
\midrule
$\Delta$ & $7$ & $\binom{4}{2}{+}1$ & exact \\
$N$ & $192$ & $|S_4|{\times}2^3$ & exact \\
$x$ & $7/192$ & $\Delta/N$ & exact \\
$\eta_R$ & $185/192$ & $1{-}x$ & Planck \\
$\alpha_0$ & $7/960$ & Thm.~\ref{thm:alpha} & 0.08\% \\
$S_0$ & $43/7$ & Thm.~\ref{thm:S0} & struct.\ \\
$H$ & $101/6$ & Thm.~\ref{thm:H} & 0.09\% \\
$\kappa_C$ & $3/(4\pi)$ & G--B & struct.\ \\
$C^*$ & $7/8$ & $\Delta/2^3$ & struct.\ \\
\bottomrule
\end{tabular}
\end{table}

%═══════════════════════════════════════════════════════════
\section{Gauge Structure from $K_4$}
\label{sec:gauge}
%═══════════════════════════════════════════════════════════

\begin{theorem}[Gauge group]
\label{thm:gauge}
$\mathrm{SU}(3){\times}\mathrm{SU}(2){\times}\mathrm{U}(1)$ from $S_4$ on $K_4$:
\begin{enumerate}[label=(\roman*),leftmargin=2em,itemsep=2pt]
\item $\mathrm{SU}(3)$: faces, $g_3^2=3/\pi$,
\item $\mathrm{SU}(2)$: edge pairs, $g_2^2=3/(2\pi)$,
\item $\mathrm{U}(1)$: vertices, $g_1^2=9/(16\pi)$.
\end{enumerate}
\end{theorem}
\textit{Proof sketch.}
Face holonomy with 3 vertices per triangle generates $\mathrm{SU}(3)$.
Three opposite-edge pairs generate $\mathrm{SU}(2)$.
Diagonal vertex-phase gives $\mathrm{U}(1)$.
Couplings: $g_k^2=\kappa_C\times(\text{factor})$.

Ratios $48:24:9$ give $\sin^2\theta_W=3/11\approx 0.273$ at unification.
All couplings carry the $1/\pi$ consumed-loop factor.

%═══════════════════════════════════════════════════════════
\section{Electroweak Symmetry Breaking}
\label{sec:ewsb}
%═══════════════════════════════════════════════════════════

\begin{theorem}[Framing vacuum]
\label{thm:ewsb}
$\sigma_i=+1$ $(\forall i)$ breaks $\mathrm{SU}(2)\to\mathrm{U}(1)_{\tau_3}$.
\end{theorem}
\textit{Proof sketch.}
Aligned framing minimizes $\eta_O(\sigma_i{-}\sigma_j)^2$; $\tau_3$ preserves the vacuum; $\tau_1,\tau_2$ flip framings and acquire mass.
$\mathrm{SU}(3)$ is unbroken (faces $\perp$ edges); confinement from Sec.~\ref{sec:massgap}.

Photon: $A_\mu=\cos\theta_W B_\mu+\sin\theta_W W^3_\mu$.
Higgs: framing-alignment amplitude mode.

%═══════════════════════════════════════════════════════════
\section{Yang--Mills Mass Gap}
\label{sec:massgap}
%═══════════════════════════════════════════════════════════

\begin{theorem}[Mass gap]
\label{thm:massgap}
$\Delta_{\mathrm{mass}}\geq 3/\pi>0$.
\end{theorem}
\textit{Proof sketch.}
The Hodge Laplacian on $K_4$ has eigenvalue $\Delta_{\mathrm{curv}}=4$ ($S_4$-degenerate).
Gap $3/\pi\approx 0.955$ survives the continuum: IFS growth $\eta_R/r\approx 1.53>1$ prevents closure.
Banach uniqueness with $\Delta/N<1-r$ (90\% margin).
The discrete curvature is computed in the Regge calculus sense~\cite{regge1961}.

%═══════════════════════════════════════════════════════════
\section{Emergent Geometry}
\label{sec:geometry}
%═══════════════════════════════════════════════════════════

\begin{theorem}[Continuum limit]
\label{thm:continuum}
$S_{\mathrm{EH}}=(16\pi G)^{-1}\!\int\! d^4x\sqrt{g}\,R$ is the unique Banach fixed point of IFS coarse-graining, with $G=\kappa_C a_0^2$.
\end{theorem}
\textit{Proof sketch.}
Level $n{+}1\to n$ decimation contracts the action space.
The plaquette sum converges to $S_{\mathrm{EH}}$ by self-similarity.
Sub-leading: four channels $(\theta,a,c,b)$ with $z_k\propto(7/192)^{p_k}$.
Rate: $|S^{(n)}-S^{(\infty)}|\leq C(7/192)^n$.

The vierbein $\tilde{e}^a_\mu$ from DEC face normals~\cite{hirani2003} yields $g_{\mu\nu}$ and Einstein equations~\cite{einstein1915} in the continuum limit; the four-channel accommodation architecture organizing departures from GR is developed in Ref.~\cite{mayes2026grav}.

%═══════════════════════════════════════════════════════════
\section{Matter Sectors}
\label{sec:matter}
%═══════════════════════════════════════════════════════════

A fermion is a $\mathbb{Z}_2$ framing defect ($\sigma_j=-1$).

\begin{theorem}[Nonlinear fermion equation]
\label{thm:fermion}
$(i\gamma^\mu D_\mu - m_{\mathrm{phys}})\psi = \lambda_{\mathrm{NL}}|\bar\psi\psi|\psi$,
with $\lambda_{\mathrm{NL}}=441/(2048\pi^2)$.
\end{theorem}
\textit{Proof sketch.}
$\gamma$-matrices from $\deg(v)=3$ hopping directions.
Integrating out the XY phase generates self-energy and cubic nonlinearity $g_{\chi\sigma}^2/J_\sigma$.

Mass hierarchy: $m_{n+1}/m_n=H=101/6$; prediction $m_\tau/m_\mu\approx 16.833$ vs.\ measured $16.818$ (0.09\%).
Three generations: fourth requires $b$-channel beyond $H^2$.
Nuclear binding follows a Weizs\"acker-like formula~\cite{weizsacker1935} with UD-derived coefficients.

%═══════════════════════════════════════════════════════════
\section{Fine-Structure Dressing}
\label{sec:dressing}
%═══════════════════════════════════════════════════════════

Leading radiative correction:
$z_\alpha = \alpha_0/\pi + O(\alpha_0^2)$,
giving $\alpha_{\mathrm{phys}}^{-1}\approx 137.0$ vs.\ measured $137.036$~\cite{planck2018}.
Each virtual loop carries holonomy $\Phi_{\mathrm{face}}=\pi$.

%═══════════════════════════════════════════════════════════
\section{Quantum Measurement and Entanglement}
\label{sec:quantum}
%═══════════════════════════════════════════════════════════

The carrier $\mathcal{C}_i=(C_\rho,C_\theta,C_\chi,C_\sigma,C_b)$ has viability $\mathcal{V}=\min_k(C_k/C_k^{\min})\geq 1$.

Five measurement laws without a measurement axiom:
\begin{enumerate}[label=(\roman*),leftmargin=2em,itemsep=1pt]
\item Branch generation: pruned by $\eta_B=7/16$.
\item Born rule from carrier viability via Gleason~\cite{gleason1957}.
\item Decoherence at rate $\eta_O N_{\mathrm{env}}$.
\item Context stability: $e^{-\nu_0 N}$ suppression.
\item Frequency bridge: viability selection.
\end{enumerate}

\begin{theorem}[Entanglement]
\label{thm:entangle}
Packets from a common splitting vertex carry correlated carrier states via $\mathbb{Z}_2$ conservation $\sigma_A\sigma_B=\sigma_{\mathrm{parent}}$.
\end{theorem}
\textit{Proof sketch.}
Carrier conservation at the splitting vertex distributes components between daughters.
Framing is multiplicative: $\sigma_A\sigma_B=\sigma_{\mathrm{parent}}$, producing a singlet or triplet.
Local transport cannot break the correlation; only decoherence ($\eta_O N_{\mathrm{env}}$) destroys it.
Born-rule weighting gives CHSH $S=2\sqrt{2}$.

Entanglement cost: $\nu_0/\pi$ per pair.
The Madelung transform~\cite{madelung1927} independently confirms the Schr\"odinger limit.

%═══════════════════════════════════════════════════════════
\section{Cosmology}
\label{sec:cosmology}
%═══════════════════════════════════════════════════════════

On large scales the emergent Friedmann equation is
\begin{equation}
\left(\frac{\dot a}{a}\right)^2 = \frac{8\pi G}{3}(\rho_m+\rho_r+\rho_\Lambda).
\label{eq:friedmann}
\end{equation}
The spectral index from coherence retention:
\begin{equation}
n_s = \eta_R = \frac{185}{192}\approx 0.9635,
\label{eq:ns}
\end{equation}
consistent with Planck $n_s=0.9649\pm 0.0042$~\cite{planck2018}.
Tensor-to-scalar ratio:
\begin{equation}
r_T = x = \frac{7}{192}\approx 0.0365,
\label{eq:rT}
\end{equation}
within LIGO/Virgo sensitivity~\cite{ligo2016}.
A GW forward model predicts four accommodation channels at $f^2,f,f^3,f^4$ scaling.

%═══════════════════════════════════════════════════════════
\section{Predictions and Falsifiers}
\label{sec:predictions}
%═══════════════════════════════════════════════════════════

\begin{table}[t]
\centering
\caption{Zero-parameter predictions.}
\label{tab:predictions}
\small
\begin{tabular}{@{}l@{\hskip 5pt}l@{\hskip 5pt}l@{\hskip 5pt}l@{}}
\toprule
Quantity & Prediction & Meas.\ & Acc.\ \\
\midrule
$\alpha_0^{-1}$ & $960/7$ & $137.04$ & 0.08\% \\
$n_s$ & $185/192$ & $0.965$ & Planck \\
$r_T$ & $7/192$ & ${<}0.036$ & consist.\ \\
$\kappa_C$ & $3/(4\pi)$ & $0.24$ & 0.5\% \\
$m_\tau/m_\mu$ & $101/6$ & $16.82$ & 0.09\% \\
$\Delta_{\mathrm{mass}}$ & ${\geq}3/\pi$ & ${>}0$ & struct.\ \\
$\sin^2\theta_W$ & $3/11$ & $0.231$ & unif.\ \\
$g_3^2{:}g_2^2{:}g_1^2$ & $48{:}24{:}9$ & --- & lattice \\
$z_\alpha$ & $\alpha_0/\pi$ & --- & leading \\
$N_{\mathrm{gen}}$ & $3$ & $3$ & exact \\
$\Phi_{\mathrm{face}}$ & $\pi$ & --- & struct.\ \\
$G$ & $\kappa_C a_0^2$ & --- & lattice \\
\bottomrule
\end{tabular}
\end{table}

Falsifiers:
(F1)~non-GR weak-field deviations,
(F2)~no controlled GR limit,
(F3)~$n_s\notin[0.960,0.968]$,
(F4)~monopoles discovered,
(F5)~$\alpha_0^{-1}\neq 960/7$,
(F6)~no stable packets in simulation,
(F7)~$r_T>0.048$.

%═══════════════════════════════════════════════════════════
\section{Computational Model}
\label{sec:computation}
%═══════════════════════════════════════════════════════════

The simulation kernel is
\begin{equation}
H_{\min} = H_\theta + H_\chi + H_\sigma + H_{\mathrm{int}},
\label{eq:Hmin}
\end{equation}
with XY rotor $H_\theta$ ($\eta_O=1/2$), transverse Ising $H_\chi$ ($J_\chi=\nu_0/6$), ferromagnetic Ising $H_\sigma$ ($J_\sigma=1/2$), and three $S_4$-invariant cross-couplings.
All couplings derived; implementable on any lattice Monte Carlo platform.
Proxy runs at $n=5{-}7$ yield $d_f\approx 1.53$.

%═══════════════════════════════════════════════════════════
\section{Discussion}
\label{sec:discussion}
%═══════════════════════════════════════════════════════════

UD shares goals with lattice gauge theory~\cite{wilson1974}, loop quantum gravity~\cite{rovelli2004}, and causal dynamical triangulations, but derives all constants from $K_4$ rather than introducing them.
The IFS~\cite{hutchinson1981,falconer2003} provides the mathematical foundation.

Compared to string theory~\cite{polyakov1987}, UD uses a 3D simplicial primitive without extra dimensions.
Compared to LQG~\cite{rovelli2004}, UD fixes the graph topology (IFS on $K_4$) rather than dynamizing spin networks, gaining tractability at the cost of full background independence.
The Madelung transform~\cite{madelung1927} provides an independent consistency check on the quantum sector.

Open computations: RG running from lattice to electroweak scale (verifying $\sin^2\theta_W$ flows from $3/11$ to $0.231$), full Yukawa dressing for the lepton tower, and numerical Jaffe--Witten spectral verification.
These are computations within the framework, not structural gaps.

A limitation: the IFS lattice is a fixed background; full background independence would require dynamical topology.
However, IFS self-similarity provides a weaker invariance: physics is identical at every level.

%═══════════════════════════════════════════════════════════
\section{Conclusion}
\label{sec:conclusion}
%═══════════════════════════════════════════════════════════

Unified Dynamics generates geometry, gauge structure, symmetry breaking, matter, entanglement, and cosmology from a single tetrahedral substrate with one input scale and zero free parameters.
The factor $1/\pi$ appears universally as the consumed-loop conversion between geometry and matter.
The mass gap, gauge group, Weinberg angle, fine structure, mass hierarchy, Born rule, entanglement, and Einstein--Hilbert action are all derived.

Open frontiers: quantitative RG running, the full mass spectrum, and simulation at $n\geq 8$.
Falsifier F6 is decisive: if no stable packets emerge in simulation, the substrate ontology fails.

\section*{Acknowledgments}
The authors thank iterative computational dialog and theory-development workflows that enabled progressive consolidation of the Unified Dynamics framework.

\section*{Author contributions}
B.W.M.\ conceived the framework, derived all theorems and predictions, and wrote the manuscript.
A.E.M., D.H., B.R.S., J.A.D., B.B., A.C.B., and J.E.\ contributed through discussion, review, and validation of the derivation chain.

\smallskip
\noindent\textbf{Data availability statement.}
This paper has no associated data.
This is a theoretical study and no experimental data were generated or analyzed.

\begin{thebibliography}{99}
\bibitem{einstein1915} A.~Einstein, ``Die Feldgleichungen der Gravitation,'' Sitzungsber.\ Preuss.\ Akad.\ Wiss., 844--847 (1915).
\bibitem{yangmills1954} C.~N.~Yang and R.~L.~Mills, ``Conservation of isotopic spin and isotopic gauge invariance,'' Phys.\ Rev.\ \textbf{96}, 191--195 (1954).
\bibitem{wilson1974} K.~G.~Wilson, ``Confinement of quarks,'' Phys.\ Rev.\ D \textbf{10}, 2445--2459 (1974).
\bibitem{rovelli2004} C.~Rovelli, \textit{Quantum Gravity} (Cambridge University Press, Cambridge, 2004).
\bibitem{hutchinson1981} J.~E.~Hutchinson, ``Fractals and self-similarity,'' Indiana Univ.\ Math.\ J.\ \textbf{30}, 713--747 (1981).
\bibitem{falconer2003} K.~Falconer, \textit{Fractal Geometry: Mathematical Foundations and Applications}, 2nd ed.\ (Wiley, Chichester, 2003).
\bibitem{hirani2003} A.~N.~Hirani, ``Discrete exterior calculus,'' Ph.D.\ thesis, California Institute of Technology (2003).
\bibitem{gleason1957} A.~M.~Gleason, ``Measures on the closed subspaces of a Hilbert space,'' J.\ Math.\ Mech.\ \textbf{6}, 885--893 (1957).
\bibitem{madelung1927} E.~Madelung, ``Quantentheorie in hydrodynamischer Form,'' Z.\ Phys.\ \textbf{40}, 322--326 (1927).
\bibitem{planck2018} Planck Collaboration, ``Planck 2018 results.\ VI.\ Cosmological parameters,'' Astron.\ Astrophys.\ \textbf{641}, A6 (2020).
\bibitem{ligo2016} LIGO Scientific and Virgo Collaborations, ``Observation of gravitational waves from a binary black hole merger,'' Phys.\ Rev.\ Lett.\ \textbf{116}, 061102 (2016).
\bibitem{weizsacker1935} C.~F.~von Weizs\"acker, ``Zur Theorie der Kernmassen,'' Z.\ Phys.\ \textbf{96}, 431--458 (1935).
\bibitem{polyakov1987} A.~M.~Polyakov, \textit{Gauge Fields and Strings} (Harwood Academic, Chur, 1987).
\bibitem{mayes2026dd} B.~W.~Mayes, ``Distinction Dynamics: A defect-based emergent field theory,'' Zenodo (2026), \href{https://doi.org/10.5281/zenodo.19582493}{doi:10.5281/zenodo.19582493}.
\bibitem{mayes2026grav} B.~W.~Mayes, ``Unified Dynamics as an accommodation-channel extension of general relativity,'' Zenodo (2026), \href{https://doi.org/10.5281/zenodo.19582493}{doi:10.5281/zenodo.19582493}.
\bibitem{regge1961} T.~Regge, ``General relativity without coordinates,'' Nuovo Cimento \textbf{19}, 558--571 (1961).
\end{thebibliography}

\end{document}
